\scanpage{11}
\mspage{\handwritten{25}}

Enfin, il nous sera commode d'avoir un critère permettant de déduire d\deleted{\unsure{u|a}}e la propriété $\mathrm P$ une propriété de platitude. Nous utilisero\deleted{j}ns :

\uline{Lemme} \deleted{\unsure{3|5}}\handwritten{4}.3. Soient $\mathrm A$ un anneau, $(\mathrm P_i)_{i\handwritten{\in}\mathrm I}$ une famille de $\mathrm A$-modules, \inserted{$\mathrm P$ leur prod\handwritten{uit}} $\mathrm M$ un $\mathrm A$-module de présentation finie. Alors l'homomorphisme canonique
\[
\mathrm M\otimes\mathrm P\handwritten{\to\prod}\mathrm M\otimes\mathrm P_i
\]
est un isomorphisme.

En effet, écrivant $\mathrm M$ comme \deleted{\unsure{un}}comme un conoyau d'un homomorphisme de modules libres \inserted{de type fini} $\mathrm L'\handwritten{\to}\mathrm L$, et notant que le foncteur produit est exact donc exact à gauche, on est ramené au cas où $\mathrm M=\mathrm A^n$, $n$ entier $\handwritten{>}0$, où c'est trivial.

\uline{Corollaire} \deleted{3}\handwritten{4}.4. Supposons $\mathrm A$ tel que tout module de type fini sur $\mathrm A$ soit de présentationfinie, par exemple $\mathrm A$ noethérien. Alors tout produit de $\mathrm A$-modules plats est un $\mathrm A$-module plat.

En effet, il suffit de montrer que $\mathrm P\otimes$ transforme un monomorphisme $\mathrm M\handwritten{\to}\mathrm N$ en un monomorphisme, quand $\mathrm M$ et $\mathrm N$ sont de type fini; le lemme 3 donne alors immédiatement le résultat voulu.

\uline{Corollaire} \deleted{3}\handwritten{4}.5. Supposons que $\mathrm A$ satisfasse \deleted{les conditions}\handwritten{la condition} de \deleted{3}\handwritten{4}.4.\inserted{par exemple $\mathrm A$ noethérien} Soit $\mathrm P$ un $\mathrm A$-module \deleted{topologique} isomorphe à \deleted{une}la limite projective \deleted{\unsure{``stricte'' dénomb}} \deleted{de $\mathrm A$-modules projectifs de type fini (``stricte'' = par un système projectif strict),} \deleted{``dénombrable'' par un système proj\unsure{e|x}} \deleted{Alors $\mathrm P$ est un $\mathrm A$} d'un système projectif strict,\deleted{de $\mathrm A$} admettant un système cofinal dénombrable, de $\mathrm A$-modules projectifs. Alors $\mathrm P$ est un $\mathrm A$-module plat.

En effet, on voit faci\deleted{\unsure{u}}lement que sous les conditions dites, $\mathrm P$ est isomorphe à un produit de $\mathrm A$-modules projectifs, donc plats, d'où le résultat.
% archive sheet number: 10
